\documentclass[11pt]{article}
\input{style-sujet}

\begin{document}

\entetesujet{Sujet bac 2012 -- Série C}

% ====================== EXERCICE 1 ======================
\exo{1}{3 points}

\begin{enumerate}
  \item Résoudre dans $\Cb$, l'équation $Z^4 - \sqrt{2}\,Z^3 - 4\sqrt{2}\,Z - 16 = 0$ sachant qu'elle admet deux solutions imaginaires pures.
  \item On considère dans le plan complexe muni d'un repère $(O,\vec{u},\vec{v})$, les points $A$, $B$, $C$ et $D$ d'affixes respectives : $z_A = 2i$ ; $z_B = -\sqrt{2}$ ; $z_C = -2i$ et $z_D = 2\sqrt{2}$.
  \begin{enumerate}[label=\alph*.]
    \item Représenter les points $A$, $B$, $C$ et $D$ dans le repère $(O,\vec{u},\vec{v})$.
    \item Montrer que les points $A$, $B$, $C$ et $D$ appartiennent à un même cercle $(C)$ dont on précisera le rayon et le centre.
  \end{enumerate}
\end{enumerate}

% ====================== EXERCICE 2 ======================
\exo{2}{5 points}

Pour tout entier naturel non nul $n$, on considère la suite $(I_n)$ définie par :
\[ I_n = \int_0^1 x^n\, e^{-\frac{x^2}{2}}\,dx \]

\begin{enumerate}
  \item Calculer $I_1$.
  \item Montrer que $\forall n \in \N^*,\ I_n \ge 0$.
  \item Par une intégration par parties, montrer que $I_n = -e^{-\frac{1}{2}} + (n-1)I_{n-2}$ (On pourra écrire $x^n = x^{n-1}\cdot x$).
  \item Étudier le sens de variation de la suite $(I_n)$. En déduire que la suite $(I_n)$ est convergente et converge vers une limite $l$.
  \item Montrer que $\forall n \in \N^*,\ 0 \le I_n \le \dfrac{1}{n+1}$. Calculer $l$.
\end{enumerate}

% ====================== PROBLÈME ======================
\prob{12 points}

Le plan $P$ est orienté. $ABCD$ est un carré de sens direct, de centre $O$. $I$, $J$, $K$ et $L$ sont les milieux respectifs des segments $[AD]$, $[AB]$, $[BC]$ et $[CD]$.

\begin{enumerate}
  \item $E$ est le point du plan tel que le triangle $IDE$ soit rectangle isocèle en $D$ et de sens direct. Montrer que $IE = AO$.
  \item En déduire qu'il existe une rotation $r$ transformant $I$ en $A$ et $E$ en $O$. Préciser une mesure $\theta$ de l'angle de la rotation $r$.
  \item Construire $\Omega$ le centre de la rotation $r$.
  \item On désigne par $\Omega_1$, le point d'intersection des droites $(IE)$ et $(OA)$. Montrer que les points $\Omega$, $E$, $O$, $\Omega_1$ sont situés sur un cercle $(C)$ que l'on tracera.
  \item Montrer que $A\Omega_1 I\Omega$ est un carré.
  \item Donner les caractéristiques de la similitude plane directe $S$ qui transforme le carré $ABCD$ en $A\Omega_1 I\Omega$.
  \item Placer $K' = S(K)$ puis $L' = S(L)$.
  \item Soit $\overline{S} = h_{(Q,\frac{1}{2})} \circ S_{OD}$, $Q$ étant un point de la droite $(OD)$. Caractériser $\overline{S}$.
  \item On se propose de construire $Q$ sachant que $\overline{S}(C) = J$.
  \begin{enumerate}[label=\alph*.]
    \item Montrer que si $\overline{S}(C) = J$ alors $\vv{QJ} = \dfrac{1}{2}\vv{QA}$.
    \item Construire alors le point $Q$.
  \end{enumerate}
  \item Soit $f$ l'application de $P$ dans $P$ qui à tout point $M$ associe $M'$ tel que $\vv{HM'} = \dfrac{1}{2}\vv{HM}$, $H$ étant le projeté orthogonal de $M$ sur la droite $(OE)$.
  \begin{enumerate}[label=\alph*.]
    \item Caractériser $f$.
    \item Tracer $(C')$ l'image de $(C)$ par $f$.
    \item Donner la nature de la courbe $(C')$.
  \end{enumerate}
\end{enumerate}

\end{document}
